\documentclass[12pt]{scrreprt}
\usepackage{listings}
\usepackage{underscore}
\usepackage{graphicx}
\usepackage[T1]{fontenc}
\usepackage[utf8]{inputenc}
\usepackage[brazilian]{babel}
\usepackage[bookmarks=true]{hyperref}
\usepackage{float} % Pacote para forçar a posição da imagem com [H]
\usepackage{fancyhdr} % Pacote para rodapés customizados
\usepackage{lastpage} % Pacote para contar o total de páginas

% ==========================================
% CONFIGURAÇÃO DO RODAPÉ (FancyHdr)
% ==========================================
\pagestyle{fancy}
\fancyhf{} % Limpa cabeçalhos e rodapés padrão
\renewcommand{\headrulewidth}{0pt} % Remove linha do cabeçalho

% Configuração para páginas normais
\fancyfoot[L]{Aluno: Gilvan Daniel da Silva \\ Documento: PDD -- Portal Fake}
\fancyfoot[R]{Página \thepage\ de \pageref{LastPage}}

% Forçar o rodapé nas páginas de início de capítulo (estilo plain)
\fancypagestyle{plain}{
    \fancyhf{}
    \renewcommand{\headrulewidth}{0pt}
    \fancyfoot[L]{Aluno: Gilvan Daniel da Silva \\ Documento: PDD -- Portal Fake}
    \fancyfoot[R]{Página \thepage\ de \pageref{LastPage}}
}

\hypersetup{
    bookmarks=false,    % show bookmarks bar?
    pdftitle={PDD - Portal Fake},    % title
    pdfauthor={Gilvan Daniel da Silva},                     % author
    pdfsubject={PDD - Portal Fake},                        % subject of the document
    colorlinks=true,       % false: boxed links; true: colored links
    linkcolor=blue,       % color of internal links
    citecolor=black,       % color of links to bibliography
    filecolor=black,        % color of file links
    urlcolor=purple,        % color of external links
    linktoc=page            % only page is linked
}
\def\myversion{1.0 }

\begin{document}

% ==========================================
% 1. CAPA
% ==========================================
\begin{titlepage}
\begin{flushright}
    \rule{16cm}{5pt}\vskip1cm
    \begin{bfseries}
        \Large{[Nome da Instituição]}\\
        \vspace{0.5cm}
        \Large{Disciplina: Técnicas de Hyper Automation}\\
        \vspace{1.5cm}
        \Huge{PDD -- Portal Fake}\\
        \vspace{1.5cm}
        \LARGE{Versão \myversion}\\
        \vspace{1.5cm}
        Aluno: Gilvan Daniel da Silva\\
        \vspace{0.5cm}
        Professor: Moisés Levy\\
        \vspace{0.5cm}
        Semana: 2\\
        \vspace{1.5cm}
        Data: \today\\
    \end{bfseries}
\end{flushright}
\end{titlepage}

% ==========================================
% 2. SUMÁRIO
% ==========================================
\tableofcontents
\clearpage

% ==========================================
% 3. CORPO DO DOCUMENTO
% ==========================================
\chapter{Introdução}

\section{Objetivo do Processo}
O Portal Fake faz registro de pessoas, armazena dados, é um CRUD básico mas funcional e de extrema importância para qualquer empresa. A empresa que o usa contém em seu corpo de colaboradores, usuários que atualmente realizam em torno de 250 cadastros por semana, visando tanto a melhoria de atividades atuais como melhorias para atividades futuras, resolvi implantar técnicas de hyper automação nesse serviço, o resultado esperado é triplicarmos a quantidade de registros mensais, diminuição nas tarefas repetitivas. O objetivo esperado é o aumento da produtividade.

\section{Escopo do Processo}
O processo de registros começa no campo Cadastrar usuário e termina em finalizar cadastro.

\section{Responsável pelo processo}
Usuários responsáveis por realizar o cadastro de pessoas.

\chapter{Mapeamento do Processo}

\section{Descrição do processo atual (AS-IS)}
O processo é manual e segue o seguinte padrão:
\begin{enumerate}
    \item O usuário acessa o Portal Fake
    \item Clica em Novo Cadastro
    \item Preenche os dados
    \item Clica em Salvar
    \item O sistema valida as informações
    \item O cadastro é armazenado
\end{enumerate}

\begin{figure}[H]
    \centering
    \includegraphics[width=\textwidth]{portal.png}
    \caption{Fluxograma do Processo Atual (AS-IS)}
    \label{fig:processo_atual}
\end{figure}

\section{Processo futuro (TO-BE)}
Após mapeamento das funcionalidades e requisitos tanto para o sucesso quanto para a falha na realização dos processos do sistema, é possível visualizar melhorias a serem implementadas, entre elas:

Automatizar o processo de cadastro desde a importação de dados, preenchimento de dados, validação de CPF, e até mesmo proteção física dos colaboradores ao evitar tarefas repetitivas, o que muitas das vezes gera passivos, afastamentos, erros humanos devido o cansaço. Todos esses benefícios através de RPA que identifica no domínio do portal fake, os ids dos botões, e os preenche com os campos de um csv.

Para o tratamento de erros, a criação de logs para registro de um cpf não que não é encontrado/inválido, e evitando assim que o processo seja interrompido ou até mesmo que se perca o registro de cpfs inválidos.

\chapter{Métricas e Regras}

\section{Volume do Processo}
Centenas de vezes.

\section{Frequência do Processo}
Como está agora é possível que semanalmente, a empresa consiga realizar em média 250 cadastros de pessoas, o que gera em torno de 1000 cadastros por mês.

\section{Regras de negócio}
\begin{itemize}
    \item \textbf{Nome:} Obrigatório
    \item \textbf{CPF:} Obrigatório ter 11 números e não deve se repetir
    \item \textbf{E-mail:} Deve ser válido
\end{itemize}

\chapter{Requisitos do Processo}

\section{Entradas do Processo}
\begin{itemize}
    \item Nome
    \item Sobrenome
    \item CPF
    \item E-mail
    \item Telefone
    \item Endereço
\end{itemize}

\section{Saídas do Processo}
Salva o registro.

\section{Exceções e Erros}
\begin{itemize}
    \item Logs de erro: CPF inválido
    \item CPF duplicado
    \item E-mail inválido
    \item Campos vazios
\end{itemize}

\chapter{Análise de Viabilidade}

\section{Viabilidade da Automação}
Após levantamento concluí-se que o projeto é viável.

\section{ROI (Retorno sobre investimento)}
\begin{itemize}
    \item Aumenta a produtividade em até 3x ou mais
    \item Diminui erro humano
    \item Possibilita ganho de tempo para executar outros processos
    \item Torna processos padronizados
    \item Reduz o retrabalho
    \item Reduz o desgaste do colaborador
    \item Reduz riscos de afastamento por DORT
\end{itemize}

\end{document}
